% Generated draft: requires mathematical and source review.
\documentclass[11pt]{article}
\usepackage[margin=25mm]{geometry}
\usepackage{fontspec}
\setmainfont{Latin Modern Roman}
\usepackage{amsmath,amssymb,mathtools,unicode-math}
\setmathfont{Latin Modern Math}
\usepackage{enumitem,needspace,xurl,hyperref}
\setlength{\parindent}{0pt}
\setlength{\parskip}{5pt}
\setlength{\emergencystretch}{3em}
\newcommand{\N}{\mathbb{N}}
\newcommand{\Z}{\mathbb{Z}}
\newcommand{\Q}{\mathbb{Q}}
\newcommand{\R}{\mathbb{R}}
\newcommand{\C}{\mathbb{C}}
\newcommand{\F}{\mathbb{F}}
\newcommand{\A}{\mathbb{A}}
\newcommand{\PP}{\mathbb{P}}
\newcommand{\E}{\mathbb{E}}
\newcommand{\eps}{\varepsilon}
\newcommand{\dd}{\,\mathrm{d}}
\begin{document}
\section*{UnsolvedMath ID: 800289}
\textbf{Source problem code:} catalogue entry
\subsection{Definitions and mathematical statement}
Let \(I\subseteq\R\) denote the nondegenerate interval intended in the source, and let \(\operatorname{Homeo}_{+}(I)\) be the group, under composition, of strictly increasing homeomorphisms \(I\to I\). The supplied statement does not specify the endpoints of \(I\) or whether they belong to \(I\); this interval convention remains unresolved.

Use the usual discrete-group interpretation of a \emph{Kazhdan group}: a group \(G\) has Kazhdan's property \((T)\) if every unitary representation \(\pi:G\to\mathcal U(H)\) on a complex Hilbert space that has almost invariant unit vectors has a nonzero invariant vector. Here almost invariant unit vectors means that, for every finite subset \(F\subseteq G\) and every \(\varepsilon>0\), there exists \(v\in H\) with \(\|v\|=1\) such that
\[
\|\pi(g)v-v\|<\varepsilon\qquad(g\in F).
\]
An invariant vector is a vector \(w\ne0\) satisfying \(\pi(g)w=w\) for every \(g\in G\). No topology inherited from the homeomorphism group is used in this definition.

\textbf{Question (Navas, 2017, Question 3).} For the interval \(I\) intended by the source, does there exist a subgroup
\[
G\leq\operatorname{Homeo}_{+}(I),\qquad G\ne\{\operatorname{id}_I\},
\]
with Kazhdan's property \((T)\)? As the source states, such a nontrivial subgroup would necessarily be infinite. Describing \(G\) as a subgroup requires a faithful action on \(I\). The interval convention must be confirmed before treating this as a fully specified statement.
\subsection{Short English statement}
Navas asks whether a nontrivial, necessarily infinite, group with Kazhdan's property \((T)\) can act faithfully by orientation-preserving homeomorphisms of an interval. The supplied record labels the question open; that status has not been independently verified. The interval's endpoint convention is unspecified.
\subsection{Sources}
\begin{itemize}
\item Local UnsolvedMath registry, permanent ID \(800289\), exact source problem code \texttt{catalogue entry}, supplied complete source record; upstream revision \texttt{d371ee917ac9647fbb06cd6fc529deec6f3dda82}. The record attributes the question to Navas (2017) and labels it \emph{open}.
\item Navas, \emph{Group Actions on 1-Manifolds: A List of Very Concrete Open Questions} (2017), Question 3 (\#3). Supplied source link: \url{https://arxiv.org/abs/1712.06462}. The record reports access on 2026-10-16 and extraction by \texttt{arxiv\_latex}; these are supplied provenance details, not a claim of independent source access or verification.
\end{itemize}
\end{document}
